\documentclass[11pt]{article}
\usepackage[margin=1in]{geometry}
\usepackage{amsmath,amssymb,amsthm,hyperref}
\hypersetup{hidelinks}
\emergencystretch=2em
\newtheorem{theorem}{Theorem}
\newcommand{\C}{\mathbb C}
\newcommand{\Z}{\mathbb Z}
\newcommand{\F}{\mathcal F}
\title{Conjecture 00000002520:\\Separated bases for finite tensor Fourier transforms}
\author{GitHub: gaochengzhecpu}
\date{}
\begin{document}
\maketitle

\begin{abstract}
For any finite number of axes of arbitrary positive lengths, the
multidimensional discrete Fourier transform is the tensor product of
the one-dimensional transforms. We prove the full forward and inverse
formulas, preservation of separability in both directions, and the
product formula for an actual Fourier basis. Lean verifies the result
for all finite shapes using Mathlib's discrete Fourier linear
equivalence, tensor products, tensor-product bases, and canonical
coordinate isomorphisms.
\end{abstract}

\section{Scope and precise statement}
We read the source's tensor Fourier assertion in the standard setting
of finite dimensional complex tensors and the discrete Fourier
transform (DFT). The English statement specifies neither a continuous
domain nor a function space. Its Chinese text repeats the terms for
tensor, high-dimensional Fourier transform, and separated basis, but
does not supply an additional coherent mathematical condition.
The precise finite discrete statement below is the interpretation proved
here. We make no claim about Fourier integrals on Euclidean spaces,
infinite tensor products, or analytic convergence questions.

Let $d\ge0$ and let $N_1,\ldots,N_d$ be arbitrary positive integers.
Put
\[
G=\prod_{i=1}^d\Z/N_i\Z,\qquad
\mathcal V=\C^G,\qquad M=\prod_{i=1}^dN_i.
\]
The empty product $G$ is a singleton, so $d=0$ gives $\mathcal V=\C$
and $M=1$. For each axis use the standard additive character
\[
\chi_i(a)=\exp(2\pi\mathrm i\,a/N_i),
\]
where any integer representative of $a\in\Z/N_i\Z$ gives the same
value. Define
\[
(\F f)(k)=\sum_{j\in G}\left(\prod_{i=1}^d\chi_i(-j_i k_i)\right)f(j).
\]
For $f_i:\Z/N_i\Z\to\C$, a separated array is
$f(j)=\prod_i f_i(j_i)$.

\begin{theorem}\label{thm:main}
For every finite shape $(N_1,\ldots,N_d)$:
\begin{enumerate}
\item $\F$ is a complex linear equivalence. Its inverse is
\[
(\F^{-1}g)(j)=\frac1M\sum_{k\in G}
  \left(\prod_{i=1}^d\chi_i(j_i k_i)\right)g(k).
\]
\item The transform factors as the tensor product of the
one-dimensional DFTs under the canonical identification
$\bigotimes_{i=1}^d\C^{\Z/N_i\Z}\simeq\mathcal V$.
This identity holds for every tensor, whether separated or not.
In particular,
\[
\F\left(\prod_i f_i(j_i)\right)(k)
=\prod_i(\F_i f_i)(k_i),
\qquad
\F^{-1}\left(\prod_i g_i(k_i)\right)(j)
=\prod_i(\F_i^{-1}g_i)(j_i).
\]
\item The vectors
\[
b_k(j)=\frac1M\prod_{i=1}^d\chi_i(j_i k_i),\qquad k\in G,
\]
form a basis of $\mathcal V$. Each vector is the product of its
one-dimensional Fourier basis vectors
$N_i^{-1}\chi_i(j_i k_i)$, and the coordinates of $f$ in this basis
are exactly $(\F f)(k)$.
\end{enumerate}
\end{theorem}

The normalization is deliberate: the forward transform is unnormalized,
and the inverse carries $1/M$. Thus $b_k=\F^{-1}\delta_k$, where
$\delta_k$ is the standard coordinate vector. Axis length one and the
empty shape are included.

\section{Proof}
\begin{proof}
The one-dimensional DFT on $\Z/N\Z$ is a complex linear equivalence
with inverse
\[
(\F_N^{-1}g)(j)=N^{-1}\sum_{k\in\Z/N\Z}\chi(jk)g(k).
\]
Apply this equivalence along each axis. These operations are linear
equivalences on the full array space, so their composition is again
one. Expanding the successive finite sums and interchanging their order
gives exactly the defining formula for $\F$. Reversing the composition
and expanding the one-dimensional inverse formulas gives the displayed
inverse, with normalization $\prod_iN_i^{-1}=1/M$.

For a separated input, the summand factors into independent terms:
\[
\left(\prod_i\chi_i(-j_i k_i)\right)\prod_i f_i(j_i)
=\prod_i\bigl(\chi_i(-j_i k_i)f_i(j_i)\bigr).
\]
Finite distributivity gives
\[
\sum_{j\in G}\prod_i\bigl(\chi_i(-j_i k_i)f_i(j_i)\bigr)
=\prod_i\sum_{j_i\in\Z/N_i\Z}\chi_i(-j_i k_i)f_i(j_i).
\]
This is the forward separability formula. The same calculation with
positive characters and the product inverse normalization proves the
inverse separability formula.

To identify the operator on actual tensor products, suppose the first
axis has length $N$ and the remaining grid is $G'$. The canonical
coordinate equivalence
\[
T:\C^{G'}\otimes_{\C}\C^{\Z/N\Z}
\longrightarrow \C^{(\Z/N\Z)\times G'}
\]
sends $g\otimes f$ to the array $f(j_1)g(j')$. The separability formula
shows
\[
T\circ(\F_{G'}\otimes\F_N)=\F_{(N,G')}\circ T
\]
on pure tensors. Pure tensors span the tensor product and both sides
are linear, so the identity holds for every tensor. Induction on the
number of axes proves the full factorization. The first-axis versus
remaining-axis order above only uses the canonical permutation of
tensor coordinates. The inverse identity follows by taking inverses.

Finally, $\F^{-1}$ sends the coordinate basis $(\delta_k)$ to a basis
$(b_k)$. Substituting $\delta_k$ into the inverse formula gives the
displayed expression for $b_k$. Factoring the expression proves its
product-basis form. Equivalently, the canonical tensor coordinate
equivalence carries the tensor product of the remaining-axis Fourier
basis and the first-axis Fourier basis onto the full Fourier basis.
Since $\F b_k=\delta_k$, the Fourier values are precisely the basis
coordinates. This proves every assertion, including the empty shape
where all transforms are the identity.
\end{proof}

\section{Formal verification and scope}
The Lean model records a shape as a list of natural numbers $n_i$ and
uses the positive axis length $N_i=n_i+1$. Every finite list of positive
lengths has this form; this encoding is also proved. Its recursive grid
is a singleton for the empty shape and
$\operatorname{ZMod}(n+1)\times G'$ when an axis is added.
The arrays $G\to\C$ are the usual coordinate spaces of finite tensors.

The transform \texttt{dftAll} is built from Mathlib's actual
\texttt{ZMod.dft} linear equivalence, currying, and pointwise linear
equivalences on remaining axes. It is not a symbolic operator postulated
to satisfy separation. The proof then derives the complete character
sum formulas for arbitrary arrays. Mathlib's \texttt{ZMod.stdAddChar}
is the standard cyclic-group Fourier character used above.

The formalization contains actual Mathlib tensor products and bases.
The canonical coordinate equivalence uses
\begin{center}
\texttt{TensorProduct.piScalarRight}.
\end{center}
The theorem
\begin{center}
\texttt{TensorFourier.tensor\_factorization}
\end{center}
proves the operator identity for every element of the genuine tensor
product, by tensor-product induction. Its inverse is proved separately.
The theorem
\begin{center}
\texttt{TensorFourier.tensor\_product\_basis}
\end{center}
identifies Mathlib's actual \texttt{Basis.tensorProduct} with the full
Fourier basis under those coordinates. The dimension is arbitrary in
all these statements; no fixed two-dimensional example substitutes
for the general claim.

The universal proposition proved by
\begin{center}
\texttt{TensorFourier.conjecture}
\end{center}
includes the forward and inverse character sums, forward and inverse
pure-factor separation, the full product-basis identity, actual
tensor-product operator factorization, and tensor-product basis
compatibility. All parts of Theorem~\ref{thm:main} in the stated finite
discrete setting are checked. No continuous Fourier-transform assertion
is claimed or left implicit.

The strict check treats warnings as errors. The principal theorem
audits contain only \texttt{propext}, \texttt{Classical.choice}, and
\texttt{Quot.sound}. There are no unproved placeholders or added axioms.
Development and self-review were performed by a delegated AI agent;
the coordinating agent reviews the work separately. No external
independent review is claimed.

\section{Source and reproduction}
The frozen bilingual statement is preserved in \texttt{SOURCE.md}.
The verification folder records its provenance and SHA-256.
The source is conjecture
\href{https://github.com/The-Last-Math-Competition/The-Last-Math-Competition/blob/main/conjectures/00000002520.md}{00000002520}
in the competition repository. Both languages were inspected, and the
scope decision is stated explicitly above.

The project pins Lean 4.19.0 and Mathlib commit
\begin{center}
\texttt{c44e0c8ee63ca166450922a373c7409c5d26b00b}.
\end{center}
From \texttt{lean/}, obtain the pinned dependencies and Mathlib cache,
then run \texttt{lake build} and
\begin{center}
\texttt{lake env lean -DwarningAsError=true Main.lean}.
\end{center}
The file ends with explicit axiom audits. Compile this paper with
Tectonic or another LaTeX engine supporting the listed packages.
\end{document}
